\section[Formal Interpretations: Transformation, Typing, and Local Readability]{Formal Interpretations: Transformation, Typing, and\texorpdfstring{\\}{ }Local Readability}\label{sec:formal}

The preceding sections have used mathematical language where it was needed, and this section collects the structure into three formal perspectives, category theory \citep{maclane1998}, type theory \citep{martinlof1984,pierce2002} and sheaf theory \citep{maclane1992}, each of which answers a different question about editions. Category theory distinguishes transformations and their compositions, and in particular asks when two ways of reaching the same page, or of travelling between the same endpoint alphabets, are equivalent. Type theory distinguishes the kinds of entity being transformed, and asks what a publishing system should be forced to state about an edition before it is allowed to build it. Sheaf theory asks whether interpretations acquired in different passages of a book remain compatible as letters change, at different rates, from one passage to the next. These are presented as distinct perspectives with explicit definitions, and a claim is made only where it is proved; where a perspective is used as an organising metaphor, that is said.

\subsection{Category-theoretic interpretation: transformation and composition}

\paragraph{Alphabets, strings and rendering.}
Let $\mathbf{Alph}$ be the category whose objects are finite alphabets and whose morphisms are functions between them. A substitution is an automorphism $\pi:\Sigma\to\Sigma$ in this category, and a reciprocal substitution is an automorphism with $\pi\circ\pi=\id_\Sigma$, an element of order at most two in the group $\mathrm{Aut}(\Sigma)$. The assignment $A\mapsto A^{*}$ of the set of finite strings, with the extension $f^{*}$ of a letter map $f$ to strings, is a functor from $\mathbf{Alph}$ to the category of monoids, the free-monoid functor. A glyph assignment is a morphism $g:A\to\GG$, where $\GG$ is regarded as an alphabet of shapes, and rendering a string is the morphism $g^{*}:A^{*}\to\GG^{*}$. The statement that rendering a substituted text is the same as rendering the original text with a reassigned font is now a statement about composites.

\begin{proposition}[Letterwise commuting]\label{prop:letterwise}
For every substitution $\pi:A\to A$, glyph assignment $g:A\to\GG$ and string $X\in A^{*}$, $g^{*}(\pi^{*}(X))=(g\circ\pi)^{*}(X)$.
\end{proposition}

\begin{proof}
Functoriality of $(\cdot)^{*}$ gives $(g\circ\pi)^{*}=g^{*}\circ\pi^{*}$, and applying both sides to $X$ gives the claim.
\end{proof}

This is the composition-level content of the equivalence stated informally in Section~\ref{sec:edition}, and its proof is a single line because letterwise rendering is a functor. The line also shows exactly where the statement fails to generalise. A shaping function $\sh_g$ of the kind used for Arabic is a function $A^{*}\to\GG^{*}$ that is not a monoid homomorphism, because the forms at the boundary of a concatenation $XY$ depend on both $X$ and $Y$, so $\sh_g(XY)\ne\sh_g(X)\,\sh_g(Y)$ in general. The category of monoids is therefore too coarse for contextual display, and one must work in a category whose morphisms are functions between sets of strings, with Proposition~\ref{prop:commute} supplying a hypothesis, preservation of joining classes, under which the square commutes anyway.

\paragraph{Presentations and equivalence of presentations.}
Define a \emph{presentation} to be a pair $(X,g)$ of a string $X\in A^{*}$ and a glyph assignment $g:A\to\GG$, and let a morphism $f:(X,g)\to(X',g')$ be a letter map $f:A\to A'$ with $f^{*}(X)=X'$ and $g'\circ f=g$. Then $g'^{*}(X')=g'^{*}f^{*}(X)=(g'\circ f)^{*}(X)=g^{*}(X)$, so the displayed page is invariant along morphisms. The two presentations of the example of Section~\ref{sec:edition}, the stored text PUB drawn with a font that draws P as B and B as P, and the stored text BUP drawn with the ordinary font, are isomorphic: the isomorphism is the substitution $\pi$ itself, which is why one may be said to \emph{be} the other seen through a change of stored representation. Compensated and uncompensated configurations are in this sense isomorphic presentations, and the stored-text principle is a recommendation about which object of an isomorphism class to store.

\paragraph{Three notions of sameness.}
Three different relations between presentations have been in use and should not be conflated. The first is \emph{equality of displays}, $g^{*}(X)=g'^{*}(X')$, which is preserved by morphisms of presentations. The second is \emph{recoverability}, the existence of a decoder $d:\GG^{*}\to A^{*}$ with $d\circ g^{*}=\id$. This holds if and only if $g$ is injective on $A$, which is to say if and only if no two letters share a glyph, and Proposition~\ref{prop:collapse} can be rephrased as the statement that a linear exchange of two letters necessarily passes through a presentation whose rendering map is not injective. The third is \emph{equivalence of reader performance}, the claim that readers do equally well on two editions in comprehension, recognition or transfer. This is not a relation definable in the categories above, because it is a relation between outcomes of experiments and not between the structures that the morphisms preserve. A successful character conversion establishes the first relation and, when the glyph assignment is injective, the second, and it establishes none of the psychological ones by itself.

\paragraph{Routes as paths, and why topology does not distinguish them.}
Consider the shape space $\GG$ with its topology. Routes are paths, and paths in a space form a category, in which objects are points, morphisms are paths, and composition is concatenation, associative up to reparametrisation. A cycle of presentation, such as Latin to hybrid to Standard Galactic and back, is the concatenation of routes into a loop. Two paths with the same endpoints are different morphisms of this category, and become identified only after passing to homotopy classes. In a vector space of shapes every two paths with the same endpoints are homotopic, so homotopy cannot distinguish routes of the kind discussed in Section~\ref{sec:drift}, and the relevant structure is the \emph{constrained} space of whole-alphabet configurations,
\[
\mathcal{C}_\theta=\bigl\{(y_x)_{x\in A}\in\GG^{A}:\ d(y_x,y_{x'})\ge\theta\ \text{for all }x\ne x'\bigr\},
\]
in which distinct letters are kept at least $\theta$ apart. The joint route of an edition is a path in $\mathcal{C}_\theta$ from the source configuration to the destination configuration, and two such routes are equivalent as routes if they are homotopic \emph{within} $\mathcal{C}_\theta$, which is a strictly finer relation than homotopy in $\GG$. For two letters and a shape space that is a plane, the space of pairs at distance at least $\theta>0$ is homotopy equivalent to a circle, so there are inequivalent ways of exchanging two letters, and an exchange followed by an exchange back traces the generating loop if the two exchanges turn the same way and a contractible loop if they turn in opposite ways. In shape spaces of three or more dimensions the corresponding spaces of distinct points are simply connected and every route is equivalent to every other, so the differences that remain are metric: whether a path with margin at least $\theta$ exists at all, which depends on the size of $\theta$ relative to the spacing of the shapes and, in a bounded region of legible forms, can fail, and how large the margin is along the path. The existence of a margin-$\theta$ route is the formal form of the claim that a drift preserving distinctness exists. Different routes between the same endpoints are thus not interchangeable, for three reasons that correspond to three equivalence relations of increasing strength: they can differ in whether they exist within $\mathcal{C}_\theta$, in their margins, and in the reader performance they produce, and of these the first two can be settled by mathematics and computation and the third only by experiment.

\subsection{Type-theoretic interpretation: typing the publishing system}

A weakness of informal descriptions of editions is that different kinds of object are all called text. A type-theoretic treatment declares the kinds and forces conversions to state their inputs and outputs. Let $\mathsf{Alphabet}$ be the type of finite alphabets and, for $A:\mathsf{Alphabet}$, let $\mathsf{Text}(A)$ be the type of finite lists over $A$. Linguistic text, encoded character sequences, glyph sequences and laid-out pages are given different types, and a conversion is a function between them: for a substitution $\pi$, $\mathsf{encode}_\pi:\mathsf{Text}(A)\to\mathsf{Enc}_\pi(A)$ and $\mathsf{decode}_\pi:\mathsf{Enc}_\pi(A)\to\mathsf{Text}(A)$, where $\mathsf{Enc}_\pi(A)$ is a copy of $\mathsf{Text}(A)$ indexed by the substitution that produced it, so that decoding with a substitution other than the one used for encoding is a type error and not a silent corruption. A rendering function $\mathsf{render}$ takes an element of $\mathsf{Enc}_\pi(A)$ and a font for $A$ to a glyph sequence and does not count as a translation; a translation is a function between texts of different languages and has no input in common with it. A cipher does not count as a proof of readability, because readability is not a type that any of these functions inhabits.

\paragraph{Dependent structure.}
Stronger contracts need dependent types. An \emph{edition} can be declared as a dependent record, a $\Sigma$-type, with fields
\begin{itemize}[leftmargin=2em,itemsep=1pt]
\item an alphabet $A:\mathsf{Alphabet}$ and a text $T:\mathsf{Text}(A)$;
\item a substitution $\pi:A\to A$ together with a proof that $\pi$ is a bijection, and, if reciprocal editions are required, a proof that $\pi\circ\pi=\id$;
\item routes $\rho_x:[0,1]\to\GG$ for each $x:A$, with endpoints as declared;
\item a schedule $\boldsymbol\alpha$, which in the exposure-driven case is a function of the text, $\boldsymbol\alpha:(T:\mathsf{Text}(A))\to\mathsf{Schedule}(T)$, so that the edition type genuinely depends on the text it is made from; and
\item a \emph{margin certificate}: a proof that for the declared routes and schedule the margin $m$ of Section~\ref{sec:drift} is at least the declared $\theta$.
\end{itemize}
The last field is a finite computation, since the margin is a minimum over the positions at which the schedule changes and the pairs of letters, and for linear and piecewise-linear routes the minimum can be computed exactly. A builder that demands the certificate will refuse to typeset an edition containing the collision of Proposition~\ref{prop:collapse}, and the repair, a detour or a keyframe, becomes a condition of construction and not an afterthought. This is what Section~\ref{sec:source} meant by checking the hypotheses at build time. That the schedule is a function of the text also gives a typed account of a point made there: when a correction changes the text, the schedule type changes with it, so an edition made from the corrected text is a different object, and a record of observations must name which.

\begin{proposition}[Round trip]\label{prop:roundtrip}
For every bijection $\pi:A\to A$ and every $X:\mathsf{Text}(A)$, $\mathsf{decode}_\pi(\mathsf{encode}_\pi(X))=X$, where $\mathsf{decode}_\pi$ applies $\pi^{-1}$ letter by letter; and if $\pi\circ\pi=\id$ then $\pi^{-1}=\pi$, so that the underlying maps on lists of $\mathsf{encode}_\pi$ and $\mathsf{decode}_\pi$ coincide once the index $\pi$ is forgotten. (As typed functions they are not equal, since they have different domains and codomains.)
\end{proposition}

\begin{proof}
By induction on the list $X$. For the empty list both sides are empty. If $X=x::X'$, then $\mathsf{decode}_\pi(\mathsf{encode}_\pi(x::X'))=\pi^{-1}(\pi(x))::\mathsf{decode}_\pi(\mathsf{encode}_\pi(X'))=x::X'$ by the hypothesis of the induction. If $\pi\circ\pi=\id$ then $\pi^{-1}=\pi$ and the two letterwise maps on lists coincide.
\end{proof}

\paragraph{Three kinds of claim.}
The proposition is a claim of the first kind: a theorem about the encoder, inhabited by a proof that is a finite derivation. The margin certificate is a claim of a second kind, decidable by computation on the data of the edition. A third kind concerns readers. A statement about readers can itself be represented in type theory, for instance as a proposition about a recorded data set of outcomes, or as a statistical statement over a model of those outcomes, so the difficulty is not that such statements are inexpressible. It is that neither the proof about the encoder nor the margin certificate inhabits a type of that kind: no derivation about the structure of an edition establishes that readers acquire its script. The statement is supported instead by \emph{evidence}, records with provenance (the participants, the arm, the measures, the version of the edition read), and the typing discipline consists in keeping that evidence attached to the edition it concerns and distinct from the theorems and certificates, so that none of the three substitutes for another. This is the type-theoretic form of the discipline of Table~\ref{tab:status}.

\subsection{Sheaf-theoretic interpretation: local readability and global convention}

\paragraph{Setting.}
A reader meets a book locally. Glyphs are recognised in one word, inferred in another passage, and met in modified form several chapters later. To make precise what it means for such local interpretations to cohere, let $X=\{1,\dots,N\}$ be the positions of the book, let $\Gamma$ be the finite set of distinct shapes that appear, each intermediate shape counted as a shape in its own right, and let $D:X\to\Gamma$ be the display, so that $D(p)$ is the shape at position $p$. For $U\subseteq X$ write $\Gamma(U)=D(U)$ for the set of shapes appearing in $U$. A \emph{key} on $U$ is a function $k:\Gamma(U)\to\Sigma$, a reader's working assignment of letters to the shapes occurring there. Let $\Key(U)$ be the set of keys on $U$, and for $V\subseteq U$ let the restriction be $k\mapsto k|_{\Gamma(V)}$. This makes $\Key$ a presheaf on the subsets of $X$, every subset being regarded as open. A \emph{window cover} is a family $\mathcal{U}=\{U_i\}$ of subsets covering $X$, for instance runs of consecutive positions that overlap. The sheaf condition asks that every family of keys $k_i\in\Key(U_i)$ that agree on overlaps, $k_i|_{\Gamma(U_i\cap U_j)}=k_j|_{\Gamma(U_i\cap U_j)}$, glue to a unique key on $\bigcup_iU_i$ restricting to each $k_i$.

Two observations are needed before the main results. The shapes of an overlap, $\Gamma(U_i\cap U_j)$, can be a proper subset of the shapes common to the two windows, $\Gamma(U_i)\cap\Gamma(U_j)$, because a shape can occur in both windows without occurring in the part they share. Agreement on overlaps is thus weaker than agreement on shared shapes. And the sheaf condition is a condition on \emph{all} compatible families of keys, not only on the true one, so it can fail even for a book in which every shape has a definite meaning.

\paragraph{When the condition holds.}
For a shape $\gamma\in\Gamma$ let $G_\gamma$ be the graph whose vertices are the windows containing $\gamma$, with an edge between two windows when their intersection contains a position at which $\gamma$ is displayed.

\begin{proposition}[Gluing criterion]\label{prop:gluing}
(i) Gluing is unique whenever it exists. (ii) If every graph $G_\gamma$ is connected, then every compatible family of keys glues. (iii) If some $G_\gamma$ is disconnected and $|\Sigma|\ge2$, there is a compatible family of keys that does not glue, so $\Key$ is not a sheaf for $\mathcal{U}$.
\end{proposition}

\begin{proof}
(i) Every shape of $\bigcup_i\Gamma(U_i)$ lies in some $\Gamma(U_i)$, and the glued key must agree there with $k_i$. (ii) Given a compatible family, define $k(\gamma)=k_i(\gamma)$ for any window $U_i$ containing $\gamma$. If two windows are joined by an edge of $G_\gamma$, their intersection contains an occurrence of $\gamma$, so compatibility gives equal values; by connectedness the value does not depend on the window chosen, so $k$ is well defined, and by construction it restricts to each $k_i$. (iii) Suppose $G_\gamma$ has a nonempty proper union $C_1$ of components, with complement $C_2$ among the windows containing $\gamma$, and choose distinct letters $\sigma_1,\sigma_2\in\Sigma$ and any fixed $\sigma_0\in\Sigma$. Let $k_i(\gamma)=\sigma_1$ for windows in $C_1$ and $\sigma_2$ for windows in $C_2$, and let every other shape and every window not containing $\gamma$ be assigned $\sigma_0$. If two windows overlap in a position displaying $\gamma$, they are adjacent in $G_\gamma$, hence in the same component and given the same value. If they do not, their overlap does not display $\gamma$, and the keys agree on all other shapes. So the family is compatible; but any glued key would need $k(\gamma)=\sigma_1$ and $k(\gamma)=\sigma_2$.
\end{proof}

The criterion has an immediate interpretation for rare letters. A shape that occurs only a few times, far apart, and never in the overlap of two windows has a disconnected graph, so that a reader's evidence about it in one passage constrains nothing about its meaning in another. For such shapes the global key is a convention held by the reader, not a consequence of local compatibility, and this is a formal counterpart of the observation in Section~\ref{sec:drift} that the rarest letters are the ones whose acquisition is least supported by the book.

\paragraph{Two kinds of failure.}
The presheaf $\Key$ over all subsets of $X$ fails the sheaf condition in a way that has nothing to do with drift. Cover $X$ by singletons: the overlaps are empty, so every family of keys is compatible, and if a shape occurs at two positions the two singleton keys may assign it different letters, so no global key restricts to both. This happens for an ordinary static script, and Proposition~\ref{prop:gluing}(iii) records it as the disconnection of the occurrence graph $G_\gamma$ for the cover chosen. The sheaf condition is therefore a property of a presheaf together with a cover, and for covers fine enough to separate occurrences it is not expected to hold. The coherence statement below is deliberately different. It concerns one specific family, the true keys, and asks whether that family extends to a global key. For a static script it does, for \emph{every} cover, since each shape has one meaning throughout. For a swap-ending drift it does not for \emph{any} cover, since the meanings differ at two positions and every position lies in some window. The first failure is an artefact of how the book is cut into windows; the second is an obstruction in the edition itself, and it is only the second that bears on whether a reader can hold one convention for the whole book.

\paragraph{True local keys and global coherence.}
Suppose each window is small enough that every shape in it has a single meaning there, and let $\kappa_U\in\Key(U)$ be the \emph{true key} of a window $U$, the assignment of each displayed shape to the letter it stands for.

\begin{proposition}[Coherence]\label{prop:coherence}
The family of true keys $(\kappa_{U_i})$ is compatible on overlaps. It glues to a key on $X$ if and only if every shape has the same meaning at all of its occurrences.
\end{proposition}

\begin{proof}
If $\gamma$ is displayed in $U_i\cap U_j$, then at that position $\gamma$ has a single meaning, and by the assumption on windows $\kappa_{U_i}(\gamma)$ and $\kappa_{U_j}(\gamma)$ are both this meaning; so the true keys agree on overlaps. If every shape has one meaning throughout, the function assigning that meaning is a key on $X$ restricting to each $\kappa_{U_i}$. Conversely, if a shape has different meanings at two positions, then taking windows containing them gives $\kappa_{U_i}(\gamma)\ne\kappa_{U_j}(\gamma)$ and no single $k$ can restrict to both.
\end{proof}

The proposition applies directly to the two endpoint designs of Section~\ref{sec:drift}. A drift to new shapes, as from Latin to Standard Galactic along routes whose intermediate shapes are distinct from the endpoints and from one another, gives each shape a single meaning throughout, and the true keys glue to a global key. A drift that ends in swapped shapes does not. In the exchange of letters $a$ and $b$ with detour routes, the pure shape $A$ means $a$ at the start of the book and $b$ at the end, while in the middle, where both letters are in intermediate forms, the pure shape $A$ does not occur at all. If the windows are shorter than the middle stretch, no window contains both an early and a late occurrence of $A$ and none of the overlaps contains $A$ at all, so the true keys of early and late windows are compatible with those of the middle windows and nevertheless have no common extension. The book is \emph{locally coherent and globally incoherent}: no two adjacent windows disagree about any shape, so the rate at which meanings change between neighbouring windows is zero, and yet the global assignment does not exist. This formalises a feature of the swap design that the margin and exposure analyses of Section~\ref{sec:drift} do not capture. Coherence is a property of the whole book, not of its neighbourhoods, and what the reader must acquire in the swap design is a key that varies with the phase of the book and that is, at the level of single windows, invisible.

\paragraph{Refinement by context.}
The incoherence can be removed by refining the shapes. Let $\Gamma'=\Gamma\times\{\text{early},\text{late}\}$ and let the display carry a phase marker, so that the pure shape $A$ in the early part is a different element of $\Gamma'$ from the pure shape $A$ in the late part. With this refinement every element of $\Gamma'$ has a single meaning and the true keys glue by Proposition~\ref{prop:coherence}. In practice the phase marker would be a cue that the reader can track: a colour, a running footer, a typographic device that changes at the transition. This is a design suggestion that follows from the formal analysis and is untested, and it amounts to making the context in which a shape is read part of what the reader learns about the shape.

\paragraph{Contextual and conventional readability.}
The key presheaf concerns conventions, assignments of letters to shapes that hold across a window. Reading by context is more flexible: a reader may decode the same shape as one letter in one word and another letter in the next, because the sentence makes the difference clear. This is captured by the presheaf $\mathrm{Tok}$ of \emph{token decodings}, in which a section over $U$ is a function $d:U\to\Sigma$ assigning a letter to each position. It is a sheaf for every cover, since functions on positions glue trivially. Let $\mathrm{Tok}_{\mathrm{adm}}(U)$ be the set of token decodings under which every complete word lying inside $U$ is a word of the lexicon, and let $\Dec(U)\subseteq\Key(U)$ be the set of keys $k$ for which $k\circ D$ lies in $\mathrm{Tok}_{\mathrm{adm}}(U)$. Restricting an admissible decoding to a smaller window keeps it admissible, since fewer words lie inside the smaller window, so both are presheaves and $\Dec$ is a subpresheaf of $\Key$. A decoding $d$ comes from a key exactly when it is constant on positions that display the same shape, $D(p)=D(q)\Rightarrow d(p)=d(q)$. Call a window \emph{contextually readable} if $\mathrm{Tok}_{\mathrm{adm}}(U)\ne\varnothing$ and \emph{conventionally readable} if $\Dec(U)\ne\varnothing$. Conventional readability implies contextual readability and the converse fails, and the failure is informative: a window in which a reading is available through context although no single assignment of letters to shapes will serve is one in which the text can be read only by decoding the same shape differently at different positions, with no stable convention available to the reader. The collision of Proposition~\ref{prop:collapse} produces exactly such windows, since a shared glyph for two letters can be read correctly only by decoding the same shape differently at different positions. The distinction is the formal counterpart, on the side of the edition, of the difference between understanding a passage through contextual assistance and holding a convention stable enough to read unfamiliar material; whether a reader holds such a convention is a further matter, taken up next.

\paragraph{Availability, acquisition and comprehension.}
The definitions above are about what an edition makes available, and three things that are easily run together need to be kept apart. A key is \emph{available} on a window $U$ if $\Dec(U)\neq\varnothing$: some assignment of letters to the shapes of $U$ turns every complete word of $U$ into a word of the lexicon. This is a property of the edition and the lexicon and says nothing about any reader. A key is \emph{acquired} on $U$ by a reader $R$ if the assignment $\kappa_R$ that the reader actually holds, a partial function from shapes to letters whose domain and content are empirical quantities, is defined on $\Gamma(U)$ and its restriction lies in $\Dec(U)$. Acquisition is a state of the reader at a time, measurable in the way recognition in isolation is measured in Section~\ref{sec:empirical}. \emph{Comprehension} is a third thing, an outcome such as accuracy on questions about the passage or reading speed, which can occur without acquisition, through contextual decoding, and can fail with it, through inattention or the difficulty of the content. Acquisition on $U$ implies availability on $U$, since the acquired key is then an admissible one, and nothing else among the three follows formally from anything else. In particular, exposure does not appear in the model as a premise from which acquisition is derived. The connection between exposure and acquisition, that a shape seen often enough comes to be held, is a hypothesis of the empirical programme and is not a theorem about editions.

\paragraph{Transfer.}
The same machinery states what transfer to a new book requires. Let the learned book have final-stage true key $k_S$, defined on the shapes $\Gamma_S$ that it displayed in final form, and let a new book display shapes $\Gamma'$ from the same script. The key glues across the two books, in the sense of Proposition~\ref{prop:gluing}, exactly on $\Gamma_S\cap\Gamma'$, and the shapes of $\Gamma'$ outside $\Gamma_S$ have no section to extend. Transfer to unfamiliar material is therefore limited, at most, to shapes that the learned book displayed in final form, a necessary condition for any reader to hold their keys on the strength of that book. How many of those keys a given reader actually acquired is the empirical question of the preceding paragraph, and by Section~\ref{sec:drift} the shapes displayed often enough to make acquisition plausible are predominantly the frequent ones.

\paragraph{The formal question.}
The perspective yields a precise form of the question with which the programme began: under what conditions can changing local representations support a coherent, transferable interpretation? Within the model the conditions on the edition are four. The displays of small windows must give each shape a single meaning (local consistency). Either no shape changes meaning across the book, as in a new-shape endpoint with collision-free routes, or the changes of meaning are indexed by a context that the reader can track (global coherence). Every window must have an admissible key available, $\Dec(U)\neq\varnothing$ (local availability). And for transfer the shapes of the target material must lie among those the learned book displayed in final form (coverage). These conditions are necessary for the formal coherence of the interpretation, and they are not sufficient for psychological acquisition: that the sheaf-theoretic structure of an edition is coherent says that acquisition is not ruled out by the structure, and the empirical programme of Section~\ref{sec:empirical} is what could show that it occurs.

\subsection{What the three perspectives jointly provide}

Table~\ref{tab:formal} sets the three perspectives side by side. Together they distinguish the composition of transformations (what produces what, and which routes to a page agree), the kinds of entity being transformed (what must be stated and checked before an edition is built), and the compatibility of interpretations across changing contexts (what a reader must hold, and what the book can and cannot ensure about it). Each leaves something unsaid. The category perspective is silent on readers, the type perspective separates the three kinds of claim without evaluating any of the third kind, and the sheaf perspective is a model of interpretation whose psychological adequacy is not established by its mathematics.

\begin{table}[ht]
\centering
\small
\renewcommand{\arraystretch}{1.25}
\begin{tabular}{>{\raggedright\arraybackslash}p{2.0cm}>{\raggedright\arraybackslash}p{3.3cm}>{\raggedright\arraybackslash}p{4.6cm}>{\raggedright\arraybackslash}p{3.0cm}}
\toprule
Perspective & Objects & Main results here & Silent on\\
\midrule
Category & alphabets, presentations, routes as paths & letterwise commuting; presentations isomorphic under substitution; routes distinguished by margin, not topology; failure of commuting under shaping & reader performance\\
Type & texts, encodings, glyphs, editions with certificates & indexed encodings; round trip; margin certificate checked at build time; three kinds of claim & whether certificates suffice for legibility\\
Sheaf & keys on windows, token decodings & gluing criterion; coherence of true keys; swap endings locally coherent but not global; contextual versus conventional readability & whether coherent keys are acquired\\
\bottomrule
\end{tabular}
\caption{The three formal perspectives.}\label{tab:formal}
\end{table}
